\documentclass[11pt,reqno]{amsart}
\usepackage[letterpaper,margin=1in]{geometry}
\usepackage{amsmath,amssymb,mathtools}
\usepackage{microtype}
\usepackage[colorlinks=true,linkcolor=blue,citecolor=blue,urlcolor=blue]{hyperref}
\hypersetup{pdftitle={Infinitely Many Exceptions to a Prime-Square Representation Problem of Erdos},
 pdfauthor={Will Howe, Grant Huffman, Jensen Kohlmeyer, Liam Kruer}}
\newtheorem{theorem}{Theorem}[section]
\newtheorem{proposition}[theorem]{Proposition}
\newtheorem{lemma}[theorem]{Lemma}
\newtheorem{corollary}[theorem]{Corollary}
\theoremstyle{remark}
\newtheorem{remark}[theorem]{Remark}
\numberwithin{equation}{section}
\newcommand{\N}{\mathbb N}
\newcommand{\Z}{\mathbb Z}
\newcommand{\R}{\mathbb R}
\newcommand{\E}{\mathbb E}
\newcommand{\Prob}{\mathbb P}
\newcommand{\ind}{\mathbf 1}
\newcommand{\ee}[1]{\mathrm e\!\left(#1\right)}
\newcommand{\normt}[1]{\left\|#1\right\|_{\R/\Z}}
\DeclareMathOperator{\lcm}{lcm}
\allowdisplaybreaks[1]
\setlength{\emergencystretch}{2em}
\title[Exceptions to an Erd\H{o}s representation problem]{Infinitely many exceptions to a prime-square representation problem of Erd\H{o}s}
\author[W. Howe]{Will Howe}
\author[G. Huffman]{Grant Huffman}
\author[J. Kohlmeyer]{Jensen Kohlmeyer}
\author[L. Kruer]{Liam Kruer}
\date{November 11, 2026}
\subjclass[2020]{Primary 11N35; Secondary 11N36, 11A07}
\keywords{Erd\H{o}s problems, prime-square residues, sieve methods, exceptional sets}
\thanks{\raggedright E-mail addresses: \href{mailto:howew@purdue.edu}{howew@purdue.edu} (W. Howe),
\href{mailto:huffma97@purdue.edu}{huffma97@purdue.edu} (G. Huffman),
\href{mailto:jenwin@purdue.edu}{jenwin@purdue.edu} (J. Kohlmeyer),
\href{mailto:kruerl@purdue.edu}{kruerl@purdue.edu} (L. Kruer).}

\begin{document}
\begin{abstract}
We prove that infinitely many positive integers have no representation
$n=ap^2+b$ with $p$ prime, $a\geq1$, and $0\leq b<p$.
More precisely, the number of such integers in $(X,2X]$ is of order
$X/\log X$. The lower bound combines a positive-proportion avoidance
estimate for large primes with first- and second-moment estimates for
a short average of the small-prime sieve. A separated-frequency
argument permits the large primes to be added in finitely many blocks.
The small-prime moments have quadratic loss in the cutoff exponent,
and a boundary estimate makes the large-prime avoidance stable under
the averaging shifts.
\end{abstract}
\maketitle

\section{The representation problem}

Erd\H{o}s asked whether every sufficiently large integer can be written as
\begin{equation}\label{eq:representation}
 n=ap^2+b,\qquad p\ \text{prime},\quad a\geq1,\quad 0\leq b<p;
\end{equation}
see \cite[p.~68, problem~3]{Erdos}. Call an integer \emph{exceptional}
if it has no representation \eqref{eq:representation}, and put
\[
 E(X)=\#\{n\in\Z:X<n\leq2X,\ n\text{ is exceptional}\}.
\]

\begin{theorem}\label{thm:main}
There are constants $c,C>0$ and $X_0$ such that
\begin{equation}\label{eq:main}
 c\frac{X}{\log X}\leq E(X)\leq C\frac{X}{\log X}
 \qquad(X\geq X_0).
\end{equation}
In particular, infinitely many positive integers have no representation
\eqref{eq:representation}.
\end{theorem}

Erd\H{o}s recorded the upper bound in \cite[p.~68]{Erdos}.
We prove the matching lower bound and recover the upper bound from
the same sieve estimate.

We first work with positive integer $X$ and write
$I=(X,2X]\cap\Z$. For a prime $p$, let
\[
 A_p(n)=\ind_{\{n\bmod p^2<p\}},\qquad
 V(z)=\prod_{p\leq z}\left(1-\frac1p\right),\qquad
 S_z(n)=\prod_{p\leq z}(1-A_p(n)).
\]
The remainder $n\bmod p^2$ is taken in $\{0,\ldots,p^2-1\}$.
Euclidean division shows that a prime $p$ gives a representation
\eqref{eq:representation} precisely when $p^2\leq n$ and $A_p(n)=1$.
The associated intervals are $[ap^2,ap^2+p-1]\cap\Z$ for $a\geq1$.

All sums and products indexed by $p$ run over primes. We use
$\mathrm e(t)=\exp(2\pi i t)$ and write $\|t\|_{\R/\Z}$ for distance
to the nearest integer. The notation $o_\alpha(1)$ means that $\alpha$
and all other exponent parameters have been fixed before $X$ tends
to infinity.

The classical prime estimates needed below are
\begin{align}
 \sum_{p\leq t}\frac1p&=\log\log t+B+o(1),&
 V(t)&\sim\frac{e^{-\gamma}}{\log t},\label{eq:mertens}\\
 \pi(t)&\ll\frac{t}{\log t},&
 \vartheta(t):=\sum_{p\leq t}\log p&\ll t.\label{eq:cheb}
\end{align}
Here $\gamma$ is Euler's constant. These are Mertens' estimates and Chebyshev's bounds
\cite{Mertens,Chebyshev}. Partial summation also gives
\begin{equation}\label{eq:logmoments}
 \sum_{p\leq t}\frac{\log p}{p}\ll\log t,
 \qquad
 \sum_{p\leq t}\frac{(\log p)^2}{p}\ll(\log t)^2.
\end{equation}

The proof uses a fixed dyadic cutoff $z=X^\alpha$.
Section~\ref{sec:large} gives a positive proportion of points avoiding
the larger primes. Section~\ref{sec:small} estimates a short average
of $S_z$ with relative error $O(\alpha^2)$ in both moments.
Section~\ref{sec:finish} combines the estimates and proves the theorem.

\section{Avoiding the large primes}\label{sec:large}

We begin with the averaging estimate that permits separated blocks of
primes to be added successively.

\begin{lemma}[Separated prime blocks]\label{lem:blocks}
Fix $0<a\leq b\leq1/4$ and $b<c\leq1/2$.
Let $\mathcal P$ be any set of primes in $[X^a,X^b]$, and set
\[
 G(n)=\sum_{p\in\mathcal P}A_p(n),\qquad
 \lambda=\sum_{p\in\mathcal P}\frac1p.
\]
Let $\mathcal Q$ be any set of primes in $[X^c,2\sqrt X]$, and let
\[
 U(n)=\ind_I(n)\prod_{q\in\mathcal Q}(1-A_q(n)).
\]
There is $\delta>0$, depending only on the fixed exponents, such that
\begin{equation}\label{eq:correlation}
 \sum_{n\in\Z}G(n)U(n)
 =\lambda\sum_{n\in\Z}U(n)+O(X^{1-\delta}).
\end{equation}
The error is uniform in $\mathcal P$ and $\mathcal Q$.
\end{lemma}

\begin{proof}
Choose $b<\eta<c$, put $H=\lfloor X^\eta\rfloor$, and fix
$0<\epsilon<\eta-b$. Take an even, nonnegative, smooth function
$\phi$, supported in $[-1,1]$ and positive near zero. For large $H$,
\[
 \nu_H(h)=\frac{\phi(h/H)}{\sum_{r\in\Z}\phi(r/H)}
 \quad(h\in\Z)
\]
is a probability measure supported on $[-H,H]$.
For each fixed positive integer $J$, finite differencing gives
\begin{equation}\label{eq:kerneldecay}
 \left|\sum_h\nu_H(h)\ee{h\theta}\right|
 \ll_J(1+H\|\theta\|_{\R/\Z})^{-J}.
\end{equation}
Indeed, the normalizing denominator is $\asymp H$, and the sum of
the absolute $J$th differences of $\phi(h/H)$ is $O_J(H^{1-J})$.
Multiplication by $(1-\mathrm e(\theta))^J$ gives
\eqref{eq:kerneldecay}, together with the trivial bound~$1$.

The Fourier expansion modulo $p^2$ is
\[
 A_p(n)=\sum_{k\bmod p^2}\widehat A_p(k)\ee{kn/p^2},
 \qquad
 \widehat A_p(k)=\frac1{p^2}\sum_{r=0}^{p-1}\ee{-kr/p^2}.
\]
In particular, $\widehat A_p(0)=1/p$ and
$|\widehat A_p(k)|\leq1/p$.
Define $G_H(n)=\sum_h\nu_H(h)G(n+h)$ and put
$D=X^\epsilon/H$. Discard from $G_H-\lambda$ the frequencies
with distance greater than $D$ from an integer. Their total absolute
contribution is
\[
 \ll_J X^{-J\epsilon}\sum_{p\leq X^b}p
 \ll_J X^{2b-J\epsilon},
\]
which is as small a power of $X$ as required after fixing $J$.

Every remaining nonzero frequency has a representative $k/p^2$ with
\[
 0<|k|\leq p^2D<p
\]
for sufficiently large $X$. It is therefore in lowest terms.
Different labels give different points of $\R/\Z$, and their
circular separation is at least $X^{-4b}\geq X^{-1}$.
There are at most $2p^2D$ retained labels for each prime, so the sum
of squared coefficients is
\begin{equation}\label{eq:mass}
 \ll\sum_{p\in\mathcal P}\frac{p^2D}{p^2}
 \ll X^bD\ll X^{b+\epsilon-\eta}.
\end{equation}

For completeness, a separated family of frequencies
$\{\theta_r\}$ with spacing at least $X^{-1}$ satisfies
\begin{equation}\label{eq:sampling}
 \sum_{n\in J_X}\left|\sum_r u_r\ee{n\theta_r}\right|^2
 \ll X\sum_r|u_r|^2
\end{equation}
on any interval $J_X$ of length $O(X)$.
To prove this, choose a nonnegative smooth majorant $w(n/X)$
for the interval after translation. Twice summing by parts gives
\[
 \left|\sum_nw(n/X)\ee{n\theta}\right|
 \ll X(1+X\|\theta\|_{\R/\Z})^{-2}.
\]
Spacing implies that each row sum of these bounds is $O(X)$.
Expanding the square and using
$2|u_ru_s|\leq|u_r|^2+|u_s|^2$ proves \eqref{eq:sampling}.
Combining \eqref{eq:mass} and \eqref{eq:sampling}, and including
the discarded frequencies, yields some $\delta_1>0$ for which
\begin{equation}\label{eq:smoothnorm}
 \sum_{X-H\leq n\leq2X+H}|G_H(n)-\lambda|^2
 \ll X^{1-\delta_1}.
\end{equation}

Enlarge and contract both the ambient interval and every forbidden
residue interval by $H$ at its endpoints. This gives indicators
$U_-$ and $U_+$ such that, for $|h|\leq H$,
\[
 U_-(n-h)\leq U(n)\leq U_+(n-h).
\]
The enlarged ambient interval lies in $[X-H,2X+H]$.
For a prime $q$, the boundary change meets at most
$O(H(X/q^2+1))$ integers in that interval. Thus
\begin{align*}
 \sum_n(U_+(n)-U_-(n))
 &\ll H+H\sum_{X^c\leq q\leq2\sqrt X}(X/q^2+1)\\
 &\ll HX^{1-c},
\end{align*}
where the reciprocal-square tail can be bounded by the sum over all
integers and $c\leq1/2$. Also $\lambda=O_{a,b}(1)$ by
\eqref{eq:mertens}.

Since $G\geq0$, averaging the pointwise comparisons and changing
variables gives
\[
 \sum_n U_-(n)G_H(n)\leq\sum_nU(n)G(n)
 \leq\sum_nU_+(n)G_H(n).
\]
Cauchy--Schwarz and \eqref{eq:smoothnorm} bound either centered
sum by $O(X^{1-\delta_1/2})$.
The boundary error is $O(HX^{1-c})$, which also saves a positive
power because $\eta<c$. This proves \eqref{eq:correlation}.
\end{proof}

Put $d=1-\log2$. We have $1/4<d<1$, since
$\exp(3/4)>1+3/4+(3/4)^2/2>2$.

\begin{proposition}\label{prop:large}
There is an absolute constant $c_0>0$ such that, for every fixed
integer $L\geq2$, with $\alpha=2^{-L}$,
\begin{equation}\label{eq:largecount}
 \#\{n\in I:A_p(n)=0\text{ for every }X^\alpha<p\leq2\sqrt X\}
 \geq c_0d^L X
\end{equation}
for all sufficiently large $X$, depending on $L$.
\end{proposition}

\begin{proof}
Each prime contributes $X/p+O(p)$ points of $I$ to its forbidden
residues. Consequently the number of points surviving all primes
$X^{1/4}<p\leq2\sqrt X$ is at least
\begin{align*}
 X-X\sum_{X^{1/4}<p\leq2\sqrt X}\frac1p
 -O\left(\sum_{p\leq2\sqrt X}p\right)
 =(d-o(1))X.
\end{align*}
Here \eqref{eq:mertens} gives the reciprocal sum $\log2+o(1)$,
and \eqref{eq:cheb} gives $\sum_{p\leq2\sqrt X}p\ll X/\log X$.

Set $K=L-2$ and $c_j=2^{-j}/4$ for $0\leq j\leq K$.
Fix $0<\tau<1/2$ and add, in the order $j=0,\ldots,K-1$, the blocks
\[
 X^{c_{j+1}}<p\leq X^{(1-\tau)c_j}.
\]
At stage $j$, the previously imposed primes are all greater than
$X^{c_j}$, so Lemma~\ref{lem:blocks} applies. The reciprocal mass
of the new block tends to
$\log(2(1-\tau))<\log2$.
If $N_j$ points survive before the block is added, the union bound
and \eqref{eq:correlation} give
\[
 N_{j+1}\geq dN_j-o(X).
\]
There are only $K$ stages, all fixed before $X\to\infty$.
The final selected-prime count is therefore at least
$(d^{K+1}-o(1))X$.

The omitted gaps are
$X^{(1-\tau)c_j}<p\leq X^{c_j}$.
Their total reciprocal mass tends to $K[-\log(1-\tau)]$.
Removing all their forbidden points costs at most
\[
 \bigl(K[-\log(1-\tau)]+o(1)\bigr)X;
\]
their combined $O(p)$ errors are $o(X)$ by \eqref{eq:cheb}.
Choose $\tau$, still at fixed $K$, so that
$K[-\log(1-\tau)]<d^{K+1}/4$.
For large $X$ at least $d^{K+1}X/2$ points remain.
This gives \eqref{eq:largecount} with $c_0=1/(2d)$.
The case $K=0$ follows directly from the initial estimate.
\end{proof}

\section{The averaged small-prime sieve}\label{sec:small}

Fix a small positive $\alpha$ and define
\begin{equation}\label{eq:average}
 z=X^\alpha,\qquad
 \eta=\alpha(1-\alpha^2),\qquad
 H=\lfloor X^\eta\rfloor,\qquad
 B(n)=\frac1H\sum_{j=0}^{H-1}S_z(n+j).
\end{equation}
The following proposition supplies the two estimates used in the
final argument.

\begin{proposition}\label{prop:moments}
There are absolute constants $C_1,C_2>0$ such that, for all
sufficiently small fixed $\alpha>0$,
\begin{align}
 \sum_{n\in I}B(n)&\geq
 (1-C_1\alpha^2-o_\alpha(1))XV(z),\label{eq:first}\\
 \sum_{n\in I}B(n)^2&\leq
 (1+C_2\alpha^2+o_\alpha(1))XV(z)^2.\label{eq:second}
\end{align}
\end{proposition}

\subsection{The first moment}

Partition the primes through $z$ into nonempty blocks
\[
 \mathcal P_j=\{p:z^{2^{-j-1}}<p\leq z^{2^{-j}}\},
 \qquad 0\leq j\leq J,
\]
ending when the lower endpoint is below $2$.
Then $J=O(1+\log\log z)$.
By \eqref{eq:mertens}, the sums $\sum_{p\in\mathcal P_j}1/p$
are bounded by an absolute constant $C_0$, uniformly in $z,j$.
Moreover,
\begin{equation}\label{eq:blocksurvival}
 V_j:=\prod_{p\in\mathcal P_j}(1-1/p)\geq v_0>0
\end{equation}
uniformly, since $-\log(1-1/p)\leq2/p$ for $p\geq2$.

Choose an absolute integer $L_0$ so large that
$C_0^k/k!\leq2^{-k}$ for $k\geq2L_0$, and put
\[
 T=\lceil\log_2(1/\alpha)\rceil+L_0,\qquad k_j=2T+2j.
\]
For $K_j(n)=\sum_{p\in\mathcal P_j}A_p(n)$, the Bonferroni
polynomials
\[
 U_j(n)=\sum_{i=0}^{k_j}(-1)^i\binom{K_j(n)}i,
 \qquad
 L_j(n)=\sum_{i=0}^{k_j-1}(-1)^i\binom{K_j(n)}i
\]
satisfy
\[
 L_j\leq\ind_{\{K_j=0\}}\leq U_j,
 \qquad U_j\geq0,\qquad U_j-L_j=\binom{K_j}{k_j}.
\]
Telescoping the product, using the nonnegativity of the upper
polynomials, gives
\begin{equation}\label{eq:bonfproduct}
 S_z(n)\geq\prod_jU_j(n)
 -\sum_j(U_j(n)-L_j(n))\prod_{i\ne j}U_i(n).
\end{equation}

Replace the variables $A_p(n)$ temporarily by independent Bernoulli
variables of means $1/p$. For a single block,
\[
 \E U_j\geq V_j,\qquad
 \E(U_j-L_j)\leq\frac{C_0^{k_j}}{k_j!}\leq2^{-k_j}.
\]
Also $\E U_j-V_j\leq\E(U_j-L_j)$, so
$\E U_j\leq V_j(1+v_0^{-1}2^{-k_j})$.
Independence between blocks and \eqref{eq:blocksurvival} show that
the expectation of the right side of \eqref{eq:bonfproduct} is at least
\[
 V(z)\left(1-O\left(\sum_j2^{-k_j}\right)\right)
 \geq V(z)(1-O(\alpha^2)).
\]

To transfer this bound to integers, expand \eqref{eq:bonfproduct}
into products of distinct $A_p$.
Every resulting squarefree modulus $m$ satisfies
\[
 m\leq z^{\sum_{j\geq0}k_j2^{-j}}
 =z^{4T+4}\leq X^{1/4}
\]
when $\alpha$ is sufficiently small.
Each coefficient has absolute value at most $J+2$.
The Chinese remainder theorem gives exactly $m$ forbidden residue
classes modulo $m^2$ for a product indexed by $m$, hence its sum
over $I$ is $X/m+O(m)$.
The total transfer error is bounded by
\[
 O\left((J+2)\sum_{m\leq X^{1/4}}m\right)
 =O((J+2)X^{1/2})=o(XV(z)).
\]
Thus $\sum_{n\in I}S_z(n)\geq(1-C_1\alpha^2-o_\alpha(1))XV(z)$.
Translating $I$ by an integer $0\leq j<H$ changes the sum of
$S_z$, which lies between zero and one, by at most $2H$.
Since $H=o(XV(z))$, averaging these translates proves
\eqref{eq:first}.

\subsection{A uniform two-position bound}\label{subsec:selberg}

For an integer $k$, let $a_p(k)$ be the cardinality of the intersection
of the forbidden residue sets at $n$ and $n+k$ modulo $p^2$.
Set
\[
 g_k(p)=\frac2p-\frac{a_p(k)}{p^2},\qquad
 V_k=\prod_{p\leq z}(1-g_k(p)).
\]
Thus $g_k(p)$ is the density of the union of these two residue sets,
and $1/p\leq g_k(p)\leq2/p$.

\begin{lemma}\label{lem:twoposition}
For sufficiently small fixed $\alpha>0$, uniformly in $k$ and in
integer intervals $J_X$ of length $X$,
\begin{equation}\label{eq:twoposition}
 \sum_{n\in J_X}S_z(n)S_z(n+k)
 \leq(1+C\alpha^2)XV_k+O(X^{3/5}),
\end{equation}
where $C$ and the implied constant are absolute.
\end{lemma}

\begin{proof}
If $g_k(p)=1$ for some $p$, then $V_k=0$ and the left side is zero.
We may therefore suppose $0<g_k(p)<1$ for every
$p\leq z$, and write $g(p)=g_k(p)$.
Extend $g$ multiplicatively to squarefree products of primes through
$z$, and put
\[
 q(r)=\prod_{p\mid r}\frac{g(p)}{1-g(p)},\qquad
 R=\lfloor X^{1/10}\rfloor,\qquad
 G_R=\sum_{r\leq R}q(r).
\]
All moduli in this proof are squarefree and supported on primes
through $z$; in particular $g(1)=q(1)=1$.

Let $F_r(n)$ indicate that the union event occurs for every $p\mid r$.
For real weights $\lambda_r$ supported on $r\leq R$, with
$\lambda_1=1$, the square
$(\sum_r\lambda_rF_r(n))^2$ majorizes $S_z(n)S_z(n+k)$.
Its main term after summation over $J_X$ is
\[
 XQ(\lambda),\qquad
 Q(\lambda)=\sum_{r,s\leq R}\lambda_r\lambda_sg([r,s]).
\]
Writing $y_t=\sum_{t\mid r}g(r)\lambda_r$, one has
\begin{equation}\label{eq:diagonal}
 Q(\lambda)=\sum_{t\leq R}\frac{y_t^2}{q(t)},
 \qquad \lambda_1=\sum_{t\leq R}\mu(t)y_t.
\end{equation}
The first identity follows by expanding
$1/g((r,s))=\sum_{t\mid(r,s)}1/q(t)$, and the second is
M\"obius inversion.
Choose $y_t=\mu(t)q(t)/G_R$.
Then $\lambda_1=1$, $Q(\lambda)=1/G_R$, and inversion gives
\begin{equation}\label{eq:weights}
 \lambda_r=\frac{\mu(r)}{G_R}
 \prod_{p\mid r}(1-g(p))^{-1}
 \sum_{\substack{s\leq R/r\\(s,r)=1}}q(s).
\end{equation}
In particular $|\lambda_r|\leq1$.
Indeed, the numerator without its sign equals
$\sum_{s\leq R/r,(s,r)=1}\sum_{e\mid r}q(se)$;
the products $se$ are distinct and all at most $R$, so this is a
subsum of $G_R$.

For $m=[r,s]$, the union conditions occupy
$\prod_{p\mid m}(2p-a_p(k))\leq m^2$ residue classes modulo $m^2$.
The Chinese remainder theorem therefore gives
\[
 \sum_{n\in J_X}F_m(n)=Xg(m)+O(m^2).
\]
There are at most $R^2$ pairs $(r,s)$ and $m\leq R^2$.
Using $|\lambda_r|\leq1$, the total error is $O(R^6)$, and hence
\begin{equation}\label{eq:selbergpre}
 \sum_{n\in J_X}S_z(n)S_z(n+k)\leq X/G_R+O(R^6).
\end{equation}

For independent Bernoulli variables $Z_p$ with means $g(p)$,
\[
 G_RV_k=\Prob\left(\sum_{p\leq z}Z_p\log p\leq\log R\right).
\]
The mean of the sum is $O(\log z)$, and its variance is
$O((\log z)^2)$ by \eqref{eq:logmoments} and $g(p)\leq2/p$.
Its second moment is therefore $O((\log z)^2)$.
Markov's inequality applied to its square gives
$G_RV_k\geq1-C'\alpha^2$ for large $X$.
After making $\alpha$ smaller if necessary,
$1/G_R\leq(1+C\alpha^2)V_k$.
Substitution in \eqref{eq:selbergpre}, with $R^6\leq X^{3/5}$,
proves the lemma.
\end{proof}

\subsection{Averaging the local factors}

For $|k|<H$, put $w_H(k)=(H-|k|)/H^2$, and put $w_H(k)=0$
otherwise. These are the probabilities of the difference of two
independent uniform integers in $\{0,\ldots,H-1\}$.
In particular, $\sum_kw_H(k)=1$ and $w_H(k)\leq1/H$.

\begin{lemma}\label{lem:localmean}
With the parameters \eqref{eq:average},
\begin{equation}\label{eq:localmean}
 \sum_kw_H(k)\frac{V_k}{V(z)^2}
 \leq(1-\alpha^2)^{-2}+o_\alpha(1).
\end{equation}
\end{lemma}

\begin{proof}
Write
\[
 R_p(k)=\frac{1-g_k(p)}{(1-1/p)^2}
 =1+\frac{a_p(k)-1}{(p-1)^2},\qquad
 v_p(k)=(R_p(k)-1)_+.
\]
The function $R_p$ has period $p^2$, is nonnegative, and is at most
$p/(p-1)$. Counting the pairs of forbidden residues gives
$\sum_{k\bmod p^2}a_p(k)=p^2$, so $R_p$ has mean one over a period.
Also
\begin{equation}\label{eq:vp}
 0\leq v_p(k)\leq\frac{a_p(k)}{(p-1)^2},\qquad
 \frac1{p^2}\sum_{k\bmod p^2}v_p(k)\leq\frac1{(p-1)^2},
 \qquad 1+v_p(k)\leq\frac p{p-1}.
\end{equation}

Choose a sufficiently small absolute $c_1>0$ and put
\[
 Y=c_1\log H,\qquad M=\prod_{p\leq Y}p^2,\qquad
 L(k)=\prod_{p\leq Y}R_p(k).
\]
By \eqref{eq:cheb}, $M\leq H^{1/4}$.
The Chinese remainder theorem shows that $L$ has mean one modulo $M$.
A uniform integer in $\{0,\ldots,H-1\}$ has residue probabilities
$1/M+O(1/H)$; convolving with another probability distribution
preserves this error bound. Consequently
\begin{equation}\label{eq:headmean}
 \sum_kw_H(k)L(k)=1+O(M/H).
\end{equation}

Set $Y_0=H^{1/4}$ and $Y_1=H^{1-\alpha^2}$.
For $Y<p\leq Y_0$, the nonnegative function $Lv_p$ has period
$Mp^2$ and, by the Chinese remainder theorem, period mean
$O(p^{-2})$. Counting complete periods and at most two incomplete
periods in $|k|<H$, and using $w_H\leq1/H$, gives
\[
 \sum_kw_H(k)L(k)v_p(k)\ll p^{-2}+M/H.
\]
By \eqref{eq:mertens} and \eqref{eq:vp},
\[
 \prod_{Y<p\leq Y_0}(1+v_p(k))\ll\frac{\log Y_0}{\log Y}.
\]
Telescoping this product and summing the preceding estimate show that
its insertion after $L$ increases the weighted mean by at most
\[
 O\left(\frac{\log Y_0}{\log Y}
       \left(\frac1Y+\frac{MY_0}{H}\right)\right)=o(1).
\]

For $Y_0<p\leq Y_1$, the same period-counting argument without $L$
gives
\[
 \sum_kw_H(k)v_p(k)\ll p^{-2}+H^{-1}.
\]
The product of $1+v_p$ in this range is $O_\alpha(1)$.
The preceding product, including $L$, is pointwise $O(\log Y_0)$
by \eqref{eq:mertens}. A second telescoping estimate therefore
bounds the additional weighted mean by
\[
 O_\alpha\left(\log Y_0
       \left(\frac1{Y_0}+\frac{Y_1}{H}\right)\right)=o_\alpha(1).
\]
Together with \eqref{eq:headmean}, this proves
\[
 \sum_kw_H(k)\prod_{p\leq Y_1}R_p(k)\leq1+o_\alpha(1).
\]
Finally,
\[
 \prod_{Y_1<p\leq z}R_p(k)
 \leq\prod_{Y_1<p\leq z}\frac p{p-1}
 =(1+o_\alpha(1))\frac{\log z}{\log Y_1}.
\]
Since $\log H\sim\eta\log X$,
\[
 \frac{\log z}{\log Y_1}\longrightarrow
 \frac{\alpha}{(1-\alpha^2)\eta}=(1-\alpha^2)^{-2}.
\]
All factors are nonnegative, so the last pointwise estimate can be
multiplied into the preceding weighted bound. This proves
\eqref{eq:localmean}.
\end{proof}

To finish Proposition~\ref{prop:moments}, expand the square in
\eqref{eq:average}. For each pair of shifts $i,j$, apply
Lemma~\ref{lem:twoposition} on $I+i$ with $k=j-i$.
Averaging the pairs and then using Lemma~\ref{lem:localmean} gives
\[
 \sum_{n\in I}B(n)^2
 \leq(1+C\alpha^2)XV(z)^2
       \bigl((1-\alpha^2)^{-2}+o_\alpha(1)\bigr)+O(X^{3/5}).
\]
For small $\alpha$, $(1-\alpha^2)^{-2}=1+O(\alpha^2)$, and
$X^{3/5}=o(XV(z)^2)$ by \eqref{eq:mertens}.
This is \eqref{eq:second}; the first moment was proved above.

\section{Stable centers and the exceptional set}\label{sec:finish}

\begin{lemma}\label{lem:stable}
Let $\alpha=2^{-L}$ with fixed $L\geq2$, and let $H$ be as in
\eqref{eq:average}. There is a set $U^*\subset I$ with
\begin{equation}\label{eq:stablecount}
 |U^*|\geq c_2d^L X
\end{equation}
for an absolute $c_2>0$ and all sufficiently large $X$, such that
every $n\in U^*$ and $0\leq j<H$ satisfy
\begin{equation}\label{eq:stableproperty}
 n+j\in I,\qquad A_p(n+j)=0
 \quad(X^\alpha<p\leq2\sqrt X).
\end{equation}
\end{lemma}

\begin{proof}
Start with the set counted in Proposition~\ref{prop:large}.
Delete its points within $H$ of either endpoint of $I$, and delete
the points meeting any forbidden residue interval enlarged by $H$
at both ends. A prime $p$ contributes at most
$O(H(X/p^2+1))$ additional deleted points. The total loss is
\[
 O\left(H+H\sum_{X^\alpha<p\leq2\sqrt X}(X/p^2+1)\right)
 \ll HX^{1-\alpha}=o_\alpha(X),
\]
because $\eta=\alpha(1-\alpha^2)<\alpha$ and $\alpha\leq1/4$.
The remaining set satisfies \eqref{eq:stableproperty}, and
\eqref{eq:stablecount} holds with $c_2=c_0/2$.
\end{proof}

\begin{proof}[Proof of Theorem~\ref{thm:main}]
We first prove the lower bound. Let $C_1,C_2$ be the constants in
Proposition~\ref{prop:moments}.
Since $d>1/4$,
\[
 \frac{(2^{-L})^2}{d^L}=(4d)^{-L}\longrightarrow0.
\]
Choose $L$ sufficiently large, put $\alpha=2^{-L}$, and set
$\rho=c_2d^L$. We can require $0<\rho\leq1$ and
$\max(C_1,C_2)\alpha^2\leq\rho/16$.
Then fix this $L$ and take $X$ sufficiently large to absorb all
$o_\alpha(1)$ errors into $\rho/16$.
Writing $V=V(z)$, we obtain
\[
 \sum_{n\in I}B(n)\geq(1-\rho/8)XV,
 \qquad
 \sum_{n\in I}B(n)^2\leq(1+\rho/8)XV^2,
 \qquad |I\setminus U^*|\leq(1-\rho)X.
\]
Cauchy--Schwarz now gives
\begin{align*}
 \sum_{n\in U^*}B(n)
 &\geq \sum_{n\in I}B(n)
   -\sqrt{|I\setminus U^*|\sum_{n\in I}B(n)^2}\\
 &\geq XV\left(1-\frac\rho8
             -\sqrt{(1-\rho)(1+\rho/8)}\right)
 \geq\frac\rho8 XV.
\end{align*}
The last inequality follows from
$(1-\rho)(1+\rho/8)\leq(1-\rho/4)^2$ for $0\leq\rho\leq1$.

Whenever $n\in U^*$ and a term $S_z(n+j)$ in $B(n)$ equals one,
the integer $m=n+j$ lies in $I$ and avoids every prime through $z$.
It also avoids the primes $z<p\leq2\sqrt X$ by
\eqref{eq:stableproperty}.
A larger prime has $p^2>2X$, so it cannot occur in a representation
of $m$. Thus $m$ is exceptional.
For any fixed $m$, at most $H$ pairs $(n,j)$ satisfy $m=n+j$ with
$0\leq j<H$. Therefore
\[
 E(X)\geq\frac1H\sum_{n\in U^*}\sum_{j=0}^{H-1}S_z(n+j)
 =\sum_{n\in U^*}B(n)\geq\frac\rho8 XV(z)
 \gg\frac{X}{\log X}
\]
by \eqref{eq:mertens}.

For the upper bound, keep one sufficiently small fixed $\alpha$.
For large $X$, every prime $p\leq z$ satisfies $p^2<X$.
Every exceptional integer in $I$ is consequently counted by $S_z$.
In Lemma~\ref{lem:twoposition}, take $k=0$.
Then $S_z(n)^2=S_z(n)$ and $V_0=V(z)$, giving
\[
 E(X)\leq\sum_{n\in I}S_z(n)
 \leq(1+C\alpha^2)XV(z)+O(X^{3/5})
 \ll\frac{X}{\log X}.
\]
This proves \eqref{eq:main} for integer $X$.
For real $X$, replacing $X$ by $\lfloor X\rfloor$ changes the
integer interval $(X,2X]$ by at most two points and changes
$X/\log X$ by a factor tending to one.
The same two bounds follow after adjusting the constants.
\end{proof}

\section*{Acknowledgment}
The authors acknowledge assistance from OpenAI's ChatGPT in developing
the argument, checking intermediate estimates, and preparing the exposition.

\begin{thebibliography}{9}
\bibitem{Erdos}
P. Erd\H{o}s, \href{https://users.renyi.hu/~p_erdos/1979-22.pdf}{\emph{Some unconventional problems in number theory}},
Mathematics Magazine \textbf{52} (1979), no.~2, 67--70.

\bibitem{Mertens}
F. Mertens, \emph{Ein Beitrag zur analytischen Zahlentheorie},
Journal f\"ur die reine und angewandte Mathematik \textbf{78} (1874), 46--62.
\href{https://doi.org/10.1515/crll.1874.78.46}{doi:10.1515/crll.1874.78.46}.

\bibitem{Chebyshev}
P. L. Tchebichef, \href{https://www.numdam.org/item/JMPA_1852_1_17__366_0/}{\emph{M\'emoire sur les nombres premiers}},
Journal de Math\'ematiques Pures et Appliqu\'ees, s\'er.~1,
\textbf{17} (1852), 366--390.
\end{thebibliography}
\end{document}
